%%%%%%%%%%%%%%%%%%%%%%%%%%%%%%%%%%%%%%%%%%%%%%%%%%%%%%%%%%%%%%%%%%%%%%
% Lớp 11 | Lý thuyết đồ thị | Đường đi Euler và Harmilton
% Dạng 1C2?2-1: Đường đi Euler
% Số câu: 3  (sinh tự động bởi tach_cau.py; từ nay thầy sửa trực tiếp file này)
%%%%%%%%%%%%%%%%%%%%%%%%%%%%%%%%%%%%%%%%%%%%%%%%%%%%%%%%%%%%%%%%%%%%%%

%%%%%%%%%%%%%%%% TRẮC NGHIỆM 4 PHƯƠNG ÁN %%%%%%%%%%%%%%%%

% nguồn: Ok_TracNghiemLop11_Lan1
\begin{ex}%[1C2Y2-1]%[ma:f29c1566]
	Trong các dãy số  sau, dãy nào là cấp số cộng?
	 \choice
		{\True $21, -3,  -27, -51, -75$}
		{$\dfrac{1}{2},  \dfrac{5}{4}, 2, \dfrac{11}{4}, \dfrac{15}{4}$}
		{$\sqrt{1}, \sqrt{2}, \sqrt{3}, \sqrt{4}, \sqrt{5}$} 
		{$\dfrac{1}{20}, \dfrac{1}{30}, \dfrac{1}{40}, \dfrac{1}{50}, \dfrac{1}{60}$} 
	 \loigiai{
		Dãy  số $21, -3, -27, -51, -75$ là cấp số cộng vì số hạng liền sau bằng số liền trước cộng với một số không đổi là $-24$.
	 }
\end{ex}


%%%%%%%%%%%%%%%% TỰ LUẬN %%%%%%%%%%%%%%%%

% nguồn: Ok_TuanLuan_Toan11_Lan1
\begin{bt}%[1C2B2-1]%[ma:8a9518b3]
	Cho cấp số cộng $\left(u_n\right)$. Tìm số  hạng đầu $u_1$, công sai $d$ trong mỗi trường hợp sau:
	\begin{multicols}{2}
		\begin{enumerate}[a)]
			\item  $u_2+u_5=42$ và $u_4+u_9=66$;
			\item $u_2+u_4=22$ và $u_1 . u_5=21$. 
		\end{enumerate}
	\end{multicols}
	\loigiai{ 
		\begin{enumerate}[a)]
			\item $u_2+u_5=42$  và $u_4+u_9=66$;\\
			Ta có $\heva{u_2+u_5=42 \\ u_4+u_9=66} \Leftrightarrow \heva{2u_1+5d=42 \\ 2u_1+11d=66}  \Leftrightarrow \heva{u_1&=11 \\ d&=4.}$\\
			Vậy số hạng đầu $u_1=11$,  công sai $d=4$
			\item $u_2+u_4=22$ và $u_1  . u_5=21$.\\
			Ta có $\heva{&u_2+u_4=22 \\ &u_1  . u_5=21} \Leftrightarrow \heva{&2u_1+4d=22 \\ &u_1.(u_1+4d)=21} \Leftrightarrow \heva{&u_1=11-2d \\ &(11-2d).(11+2d)=21}$\\
			$ \Leftrightarrow \heva{&u_1=11-2d \\ &4d^2=100}  \Leftrightarrow \hoac{&\heva{u_1=1 \\ d= 5}\\&\heva{u_1=21 \\ d= -5.}}$\\
			 Vậy có hai cấp số cộng thỏa mãn đề bài số hạng đầu $u_1=1$, công sai $d=5$; số hạng đầu $u_1=21$, công sai $d=-5$.
		\end{enumerate}
	} 
\end{bt}

% nguồn: Ok_TuanLuan_Toan11_Lan2
\begin{ex}%[1C2B2-1]%[ma:75f80d04]
	Cho cấp số cộng $\left(u_n\right)$. Tìm số hạng đầu $u_1$, công sai $d$ trong  mỗi trường hợp sau:
	\begin{multicols}{2}
		\begin{enumerate}[a)]
			\item $u_2+u_5=42$ và $u_4+u_9=66$; 
			\item $u_2+u_4=22$ và  $u_1 . u_5=21$.
		\end{enumerate}
	\end{multicols}
	\loigiai{ 
		\begin{enumerate}[a)]
			\item $u_2+u_5=42$  và $u_4+u_9=66$;\\
			Ta có $\heva{u_2+u_5=42 \\ u_4+u_9=66} \Leftrightarrow \heva{2u_1+5d=42 \\ 2u_1+11d=66} \Leftrightarrow \heva{u_1&=11 \\ d&=4.}$\\ 
			Vậy số hạng đầu $u_1=11$, công sai  $d=4$
			 \item $u_2+u_4=22$ và $u_1 . u_5=21$.\\
			Ta có $\heva{&u_2+u_4=22 \\ &u_1 . u_5=21} \Leftrightarrow \heva{&2u_1+4d=22 \\ &u_1.(u_1+4d)=21} \Leftrightarrow  \heva{&u_1=11-2d \\ &(11-2d).(11+2d)=21}$\\
			$ \Leftrightarrow \heva{&u_1=11-2d \\ &4d^2=100} \Leftrightarrow \hoac{&\heva{u_1=1 \\  d= 5}\\&\heva{u_1=21 \\ d= -5.}}$\\
			Vậy có hai cấp số cộng thỏa mãn đề bài số hạng đầu  $u_1=1$, công sai $d=5$; số hạng đầu $u_1=21$, công sai $d=-5$.
		\end{enumerate}
	} 
\end{ex}
